\documentclass[11pt]{article}
\newcommand{\DocShort}{T\'erminos de referencia}
\usepackage{fontspec}
\usepackage[spanish,es-noshorthands]{babel}
\decimalpoint
\usepackage{geometry}
\geometry{letterpaper, top=2.5cm, bottom=2.5cm, left=2.8cm, right=2.8cm}
\usepackage[expansion=false]{microtype}
\usepackage{parskip}
\setlength{\parskip}{6pt}
\usepackage{amsmath,amssymb}
\usepackage{booktabs}
\usepackage{array}
\usepackage{longtable}
\usepackage{caption}
\usepackage[dvipsnames,table]{xcolor}
\usepackage{enumitem}
\setlist[itemize]{topsep=2pt, itemsep=2pt, parsep=0pt}
\setlist[enumerate]{topsep=2pt, itemsep=2pt, parsep=0pt}
\usepackage{fancyhdr}
\usepackage[hidelinks,pdfencoding=auto]{hyperref}
\usepackage{titlesec}
\usepackage{tcolorbox}
\tcbuselibrary{breakable}
\usepackage{pdfpages}

\definecolor{darkblue}{HTML}{1B3A5C}
\definecolor{medblue}{HTML}{2E6DA4}
\definecolor{lightblue}{HTML}{D6E8F7}
\definecolor{teal}{HTML}{1A7A6E}
\definecolor{lightteal}{HTML}{D0EDEA}
\definecolor{lightgrey}{HTML}{F5F5F5}
\definecolor{midgrey}{HTML}{6F6F6F}
\definecolor{warnyellow}{HTML}{FFF3CD}

\titleformat{\section}
{\large\bfseries\color{darkblue}}
{\thesection}{0.75em}{}[{\color{medblue}\titlerule[0.8pt]}]
\titleformat{\subsection}
{\normalsize\bfseries\color{darkblue}}
{\thesubsection}{0.75em}{}
\titleformat{\subsubsection}
{\normalsize\bfseries\color{medblue}}
{\thesubsubsection}{0.75em}{}

\pagestyle{fancy}
\fancyhf{}
\fancyhead[L]{\small\color{midgrey}ICV 442 Acueductos y Alcantarillados}
\fancyhead[R]{\small\color{midgrey}\DocShort}
\fancyfoot[C]{\small\color{midgrey}\thepage}
\renewcommand{\headrulewidth}{0.4pt}
\renewcommand{\footrulewidth}{0pt}

\rowcolors{2}{lightgrey}{white}
\renewcommand{\arraystretch}{1.18}

\newtcolorbox{infobox}{
  breakable,
  colback=lightblue,
  colframe=medblue,
  boxrule=0.8pt,
  arc=3pt,
  left=8pt,
  right=8pt,
  top=7pt,
  bottom=7pt
}

\newtcolorbox{warnbox}{
  breakable,
  colback=warnyellow,
  colframe=orange!70!black,
  boxrule=0.8pt,
  arc=3pt,
  left=8pt,
  right=8pt,
  top=7pt,
  bottom=7pt
}

\newcommand{\doccover}[3]{
  \begin{center}
  {\small\color{midgrey}Pontificia Universidad Cat\'olica Madre y Maestra\\Escuela de Ingenier\'ia Civil y Ambiental}

  \vspace{1.0cm}
  {\color{medblue}\rule{\linewidth}{2pt}}
  \vspace{0.3cm}

  {\Large\bfseries\color{darkblue}ICV-442-T Acueductos y Alcantarillado}\\[0.35em]
  {\LARGE\bfseries\color{darkblue}#1}\\[0.25em]
  {\large\itshape #2}

  \vspace{0.3cm}
  {\color{medblue}\rule{\linewidth}{2pt}}
  \vspace{0.9cm}
  \end{center}

  \begin{center}
  \begin{tabular}{ll}
  \toprule
  \textbf{Asignatura} & ICV-442-T Acueductos y Alcantarillado \\
  \textbf{Profesor} & Ricardo Rafael Hern\'andez Moreira, PhD \\
  \textbf{Periodo acad\'emico} & Mayo--agosto de 2026 \\
  \textbf{Proyecto} & Dise\~no de sistemas de saneamiento para Las Charcas, Azua \\
  \textbf{Versi\'on} & #3 \\
  \bottomrule
  \end{tabular}
  \end{center}

  \vspace{0.4cm}
}


\begin{document}
\doccover{T\'erminos de referencia}{Proyecto de curso: sistemas de saneamiento para Las Charcas, Azua}{18 de mayo de 2026}

\section{Prop\'osito del proyecto}
El proyecto consiste en desarrollar, en equipos peque\~nos, un dise\~no acad\'emico de sistemas de agua potable, alcantarillado sanitario y drenaje pluvial para la comunidad de Las Charcas, provincia Azua. El producto esperado es un dise\~no detallado con memoria t\'ecnica, planos y modelaci\'on hidr\'aulica suficiente para sustentar decisiones de trazado y dimensionamiento.

\section{Competencias que se eval\'uan}
\begin{center}
\begin{tabular}{p{0.16\linewidth} p{0.68\linewidth}}
\toprule
\textbf{Competencia} & \textbf{Aplicaci\'on en el proyecto} \\
\midrule
EA3 & Obtenci\'on y uso de informaci\'on b\'asica: poblaci\'on, topograf\'ia, lluvia, normativa y brechas de datos. \\
EA4 & Aplicaci\'on de fundamentos de hidr\'aulica y justificaci\'on normativa en los tres sistemas. \\
EA5 & Consideraci\'on de factibilidad operativa, mantenimiento y coherencia general del dise\~no. \\
\bottomrule
\end{tabular}
\end{center}

\section{Alcance}
\subsection{Sistemas incluidos}
\begin{itemize}
\item \textbf{Agua potable}: fuente, conducci\'on, almacenamiento, red de distribuci\'on y accesorios principales.
\item \textbf{Alcantarillado sanitario}: colectores, pozos de visita, descarga final y bombeo solo si la topograf\'ia lo exige.
\item \textbf{Drenaje pluvial}: captaci\'on superficial, colectores, estructuras de disipaci\'on y punto de descarga.
\end{itemize}

\subsection{Nivel de detalle esperado}
\begin{itemize}
\item Memorias de c\'alculo trazables.
\item Planos legibles en escala adecuada.
\item Especificaciones t\'ecnicas de materiales y criterios de dise\~no.
\item Modelos en EPANET y SWMM con resultados verificables.
\end{itemize}

\section{Marco normativo}
\begin{enumerate}
\item Reglamento y criterios t\'ecnicos de INAPA.
\item Reglamento del MOPC cuando establezca una exigencia m\'as restrictiva o pertinente.
\item Referencias internacionales solo como apoyo o para completar vac\'ios, nunca para desplazar sin justificaci\'on el marco local.
\end{enumerate}

\begin{warnbox}
No basta citar una instituci\'on de forma gen\'erica. Las decisiones de dise\~no deben vincularse, cuando sea posible, con art\'iculos, tablas o criterios identificables de la norma consultada.
\end{warnbox}

\section{Informaci\'on disponible}
\begin{center}
\begin{tabular}{p{0.29\linewidth} p{0.18\linewidth} p{0.37\linewidth}}
\toprule
\textbf{Dato} & \textbf{Estado} & \textbf{Observaci\'on} \\
\midrule
Modelo digital de terreno & Disponible & Malla de 8 m por 8 m suministrada para el curso \\
Datos poblacionales de la ONE & Disponible & Base para proyecci\'on poblacional \\
Series de precipitaci\'on & Parcial & Pueden solicitarse a INDOMET por v\'ia institucional \\
Curvas IDF & Disponible & Ap\'endice del reglamento de dise\~no de INAPA \\
Informaci\'on sobre fuentes existentes & Incierta & Puede investigarse con INAPA \\
Usos de suelo y densidades & Parcial & Requieren interpretaci\'on a partir de cartograf\'ia e im\'agenes \\
\bottomrule
\end{tabular}
\end{center}

\section{Organizaci\'on de equipos}
\begin{itemize}
\item Los equipos estar\'an formados por 2 o 3 estudiantes.
\item Todos los integrantes deben participar en los tres sistemas.
\item Se recomienda asignar roles de coordinaci\'on, modelaci\'on, documentaci\'on y planos, con posibilidad de rotaci\'on.
\item Cada equipo debe mantener una bit\'acora semanal breve y \'util.
\end{itemize}

\section{Entregables}
\begin{center}
\begin{tabular}{p{0.33\linewidth} p{0.17\linewidth} p{0.30\linewidth}}
\toprule
\textbf{Entregable} & \textbf{Formato} & \textbf{Observaciones} \\
\midrule
Memoria de proyecci\'on poblacional & PDF + hoja de c\'alculo & Entregable independiente \\
Memoria t\'ecnica integrada & PDF & Incluye agua potable, sanitario y pluvial \\
Planos del proyecto & PDF & Planta, perfiles y detalles principales \\
Modelos hidr\'aulicos & Archivos nativos + resumen PDF & EPANET y SWMM \\
Presentaci\'on oral & PDF exportado & El archivo fuente puede elaborarse en Canva o programa equivalente \\
Anexos & PDF y hojas de c\'alculo & Patrones, bit\'acoras, tabla maestra y soportes \\
\bottomrule
\end{tabular}
\end{center}

\section{Convenci\'on de nombres para entregas}
\begin{infobox}
\texttt{ICV442\_T[Equipo]\_[Sistema]\_[Entregable]\_[Apellidos]\_v[version].pdf}
\end{infobox}

Ejemplos de uso:
\begin{itemize}
\item \texttt{ICV442\_T01\_AguaPotable\_MemoriaCalc\_HernandezPerez\_v1.0.pdf}
\item \texttt{ICV442\_T02\_Sanitario\_PlanoPerfil\_GomezRuiz\_v1.0.pdf}
\item \texttt{ICV442\_T03\_PresentacionOral\_TorresLopez\_v1.0.pdf}
\end{itemize}

\section{Coordinaci\'on con instituciones}
Si el equipo decide gestionar informaci\'on adicional, puede hacerlo mediante solicitudes formales de acceso a la informaci\'on dirigidas a las oficinas correspondientes. Esa gesti\'on es recomendable, pero no obligatoria. Si no se obtiene respuesta dentro del plazo del curso, el equipo debe documentar el intento y continuar con supuestos t\'ecnicos razonables y expl\'icitos.

\section{Criterios generales de aceptaci\'on}
\begin{itemize}
\item El trabajo debe ser coherente entre sistemas, no una suma de tres ejercicios desconectados.
\item Los supuestos deben estar declarados y ser compatibles con la informaci\'on disponible.
\item La presentaci\'on documental debe permitir que un tercero entienda qu\'e se hizo, con qu\'e datos y bajo qu\'e criterios.
\end{itemize}

\end{document}
